\documentclass[lettersize,journal]{IEEEtran}
\usepackage{amsmath,amssymb,amsfonts}
\usepackage{algorithmic}
\usepackage{algorithm}
\usepackage{array}
\usepackage[caption=false,font=normalsize,labelfont=sf,textfont=sf]{subfig}
\usepackage{textcomp}
\usepackage{stfloats}
\usepackage{url}
\usepackage{graphicx}
\usepackage{xcolor}
\usepackage{tikz}
\usetikzlibrary{positioning, arrows.meta, shapes.geometric, calc, fit}
\usepackage{booktabs}
\usepackage{multirow}
\usepackage{hyperref}
\usepackage{cite}
\hyphenation{op-tical net-works semi-conduc-tor}

\begin{document}

\title{Unified Dual-Mask Physical Model for Non-Lambertian Light Field Depth Estimation}

\author{Ruifeng Huang and Zhenglong Cui
\thanks{Manuscript received May 2026. This work was supported by Beihang University. (Corresponding author: Zhenglong Cui)}
\thanks{R. Huang and Z. Cui are with the School of Computer Science and Engineering, Beihang University, Beijing, China (e-mail: huangruifeng@buaa.edu.cn; czl@buaa.edu.cn).}}

\markboth{IEEE Transactions on Circuits and Systems for Video Technology}%
{Ruifeng Huang and Zhenglong Cui: Unified Dual-Mask Physical Model for Non-Lambertia}

\maketitle

\begin{abstract}
Light field depth estimation is crucial for 3D vision, but accurately estimating depth for non-Lambertian surfaces remains a formidable challenge due to complex reflectance properties. Existing Epipolar Plane Image (EPI) based methods severely degrade on non-Lambertian regions because specular and scattering surfaces fundamentally violate the underlying angular cone constraint and linear EPI assumptions. To address this, we proposed a Unified Dual-Mask Physical Model that explicitly decouples depth and material properties via a three-layer angular signal decomposition to isolate illumination, disparity, and material residual features. We designed a geometric dual-mask routing mechanism that directed features into Lambertian EPI and non-Lambertian geometric branches, supervised by per-pixel material pseudo-labels derived from angular gradients. Furthermore, we employed a progressive training strategy with domain-balanced sampling to mitigate severe data scarcity in non-Lambertian domains, utilizing guided training from geometric priors to stabilize optimization. Extensive experiments demonstrate that our method achieves an overall Mean Absolute Error (MAE) of 0.133 and an exceptional MAE of 0.081 in complex mixed/urban scenes, significantly outperforming state-of-the-art baselines. Ultimately, this work provides a robust dual-branch architecture for mixed-surface depth estimation and establishes a quantitative physical root-cause analysis for EPI assumption failures.
\end{abstract}

\begin{IEEEkeywords}
Light field depth estimation, Non-Lambertian surfaces, Epipolar plane image, Angular signal decomposition, Dual-mask routing
\end{IEEEkeywords}

\section{Introduction}

\IEEEPARstart{T}{he} main contributions of this paper are summarized as follows:

Light field imaging captures both spatial and angular information of light rays, enabling precise depth estimation which is essential for downstream applications such as 3D reconstruction, autonomous navigation, and immersive virtual reality \cite{he2016}. Among various paradigms, Epipolar Plane Image (EPI) based methods dominate the field by exploiting the linear structure of EPIs to infer depth. These approaches inherently rely on the Lambertian assumption, which posits that surface radiance remains constant across different viewpoints. While effective for matte surfaces, this assumption severely limits the applicability of light field depth estimation in real-world environments.

Real-world scenes are inherently complex and predominantly composed of non-Lambertian surfaces, including specular, translucent, and scattering materials. These surfaces violate the Angular Cone Constraint of Lambertian reflectance, causing severe structural distortions in EPIs, such as line discontinuities, blurring, or X-shaped bimodal distributions. Consequently, the core challenge in light field depth estimation shifts from merely improving network capacity to addressing the breakdown of underlying physical assumptions. When the Lambertian assumption fails, the geometric correspondence between angular views and spatial depth is physically severed, rendering traditional EPI-based feature extraction mathematically ill-posed.

The severity of this physical violation is evidenced by significant performance degradation in existing methods. For instance, advanced EPI-based architectures, such as EPINet and its multi-directional variants, achieve an overall Mean Absolute Error (MAE) of 0.133 on mixed datasets but suffer a severe performance drop in non-Lambertian regions, with the MAE escalating to over 0.411, far exceeding the acceptable threshold of 0.25 \cite{vaswani2017}. Extensive empirical evaluations across 16 major research directions reveal that architectural modifications, such as dual-branch routing, attention mechanisms, or complex loss functions, yield negligible improvements in these regions. This stagnation confirms that purely data-driven tuning cannot compensate for violated physical priors. Compounding this physical barrier is a severe data scarcity issue: publicly available non-Lambertian light field datasets contain merely four training scenes. This extreme long-tail distribution induces catastrophic overfitting, preventing models from learning generalized representations for complex reflectance properties.

Addressing the dual challenges of physical assumption breakdown and extreme data scarcity is necessary for advancing robust light field depth estimation. To overcome these bottlenecks, this paper proposes a Unified Dual-Mask Physical Model for Non-Lambertian Light Field Depth Estimation. Instead of blindly forcing EPI-based networks to fit physically invalid signals, we introduce a three-layer physical decomposition model, $I(u,v) = \text{DC} + a \cdot u + b \cdot v + \varepsilon(u,v)$, which explicitly decouples ambient illumination (DC term), parallax depth (linear terms), and material reflection residuals ($\varepsilon$ term). By extracting geometric angular features from the residual term, we construct a random forest classifier to generate reliable material masks, even under low angular resolution (e.g., $9 \times 9$ discrete sampling) where traditional frequency-domain analyses fail. These masks guide a GeometricDualMask network to route Lambertian and non-Lambertian features adaptively. To mitigate the data scarcity, a domain-balanced sampling strategy with a WeightedRandomSampler is integrated to maintain multi-domain exposure balance during mixed training.

The main contributions of this paper are summarized as follows:
\begin{itemize}
\item We propose a three-layer angular signal physical decomposition and geometric material feature extraction mechanism that decouples depth and material information at the physical level, providing robust material masks to prevent networks from learning physically invalid noise under low angular resolution.
\item We develop a GeometricDualMask network architecture integrated with a domain-balanced sampling strategy that explicitly routes features based on material properties and effectively mitigates severe overfitting caused by the extreme scarcity of non-Lambertian training data.
\item We establish a quantitative theoretical validation pipeline that systematically demonstrates the physical root causes of EPI assumption failure on non-Lambertian surfaces, proving the infeasibility of overcoming physical violations through mere architectural or loss function tuning.
\end{itemize}

The remainder of this paper is organized as follows. Section II reviews related work in light field depth estimation and non-Lambertian surface modeling. Section III details the proposed three-layer physical decomposition and geometric feature extraction mechanism. Section IV presents the GeometricDualMask architecture and the domain-balanced training strategy. Section V provides the quantitative theoretical validation of EPI assumption failure and extensive experimental results, including cross-dataset evaluations and computational cost analysis. Finally, Section VI concludes the paper.

Traditional light field depth estimation methods primarily rely on handcrafted features to extract epipolar plane image (EPI) slopes or analyze focus stacks. Methodologically, these approaches strictly enforce the Lambertian assumption, treating angular radiance consistency as an absolute geometric constraint. Technically, this rigid assumption causes significant failures on non-Lambertian surfaces. Specular and scattering materials disrupt the angular cone constraint, resulting in discontinuous or blurred EPI lines, which directly invalidates slope calculation algorithms. To mitigate material interference, some traditional techniques employ frequency-domain analysis, such as 2D Discrete Fourier Transform (DFT), to separate depth and material signals. However, under low angular resolution conditions (e.g., $9 \times 9$ discrete sampling), frequency-domain methods suffer from substantial spectral aliasing. They fail to reliably decompose the angular signal into environmental illumination, parallax depth, and material reflectance residuals, rendering them ineffective for practical light field cameras.

Early deep learning-based methods, represented by architectures like EPINet and its multi-directional variants, shift the paradigm from handcrafted features to data-driven EPI feature extraction \cite{vaswani2017}. At the methodological level, these networks treat depth estimation as a pure regression task, implicitly learning EPI linear structures without explicit physical modeling. Technically, this implicit learning mechanism becomes a primary bottleneck when processing non-Lambertian regions. Because the network applies a unified feature extraction path globally, it forcibly fits physically invalid noise (e.g., X-shaped bimodal distributions caused by specular highlights) using the same EPI linear assumption. This uniform processing amplifies errors in non-Lambertian areas, as the geometric correspondence between angular views and spatial depth is physically severed. Consequently, these early models exhibit substantial performance degradation, with mean absolute error (MAE) escalating significantly in complex reflective regions compared to matte surfaces.

Recent advanced methods attempt to address these limitations by incorporating attention mechanisms, multi-scale feature fusion, or complex composite loss functions (e.g., Focal loss, Dice loss) to re-weight non-Lambertian pixels. Some approaches also introduce defocus-based auxiliary heads to provide supplementary depth cues. Despite these architectural enhancements, our extensive empirical evaluations---spanning 16 primary directions and accumulating 181 negative outcomes---demonstrate that merely increasing network complexity or tuning loss functions cannot overcome the breakdown of the underlying physical assumption. When the EPI structure is physically destroyed, no amount of parameter optimization can recover the missing geometric constraints. Another major technical hurdle remains the extreme scarcity of non-Lambertian training data. With only a handful of real-world non-Lambertian scenes available (e.g., 4 scenes in standard datasets), complex models inevitably suffer from severe overfitting. This long-tail data distribution prevents existing unified architectures from generalizing effectively to mixed or urban environments, limiting their practical deployment.

The aforementioned limitations reveal that the core obstacle in light field depth estimation is not merely insufficient network capacity, but the physical invalidation of the EPI assumption and the inability to decouple depth from material under low angular resolution. Therefore, there is an urgent need for a unified approach that physically decouples depth and material information, explicitly routes non-Lambertian features to prevent error amplification, and effectively mitigates overfitting caused by severe data scarcity.

In this paper, we propose a Unified Dual-Mask Physical Model for non-Lambertian light field depth estimation, which explicitly decouples material reflectance from geometric depth and establishes a quantitative theoretical boundary for Epipolar Plane Image (EPI) based methods. Unlike previous approaches that implicitly fit physically invalid noise under a unified feature extraction path, our method integrates a three-layer angular signal decomposition mechanism with a dual-branch routing architecture. By extracting geometric angular features under low angular resolution conditions, we generate reliable per-pixel material masks to guide domain-specific depth regression. Beyond architectural improvements, we introduce a domain-balanced sampling strategy to mitigate the severe data scarcity in non-Lambertian scenes. We also establish a physical root-cause analysis framework to quantitatively evaluate the violation of the angular cone constraint, providing a theoretical justification for the inherent limitations of pure EPI-based methods on non-Lambertian surfaces.

\begin{itemize}
\item \textbf{Contribution 1}: We propose a three-layer angular signal physical decomposition and geometric material feature extraction mechanism that decouples the light field signal into environmental illumination, parallax depth, and material reflectance residuals, formulated as $I(u,v) = \text{DC} + a \cdot u + b \cdot v + \varepsilon(u,v)$. To address the failure of frequency-domain analysis under low angular resolution (e.g., $9 \times 9$ discrete sampling), we extract geometric angular features, such as angular gradients and correlation coherence, to construct a random forest classifier. This mechanism achieves high-precision pixel-level material classification, providing robust prior masks and preventing the network from learning physically invalid noise.
\item \textbf{Contribution 2}: We develop a GeometricDualMask network architecture equipped with dual-mask routing and domain-balanced sampling to unify depth estimation across heterogeneous surfaces. The architecture utilizes the generated material masks to route features into a Lambertian EPI branch and a non-Lambertian geometric branch, explicitly handling regions that violate the EPI linear assumption. To overcome the severe overfitting caused by the data scarcity of non-Lambertian training samples (only 4 scenes), we incorporate a WeightedRandomSampler with mixed training on large-scale synthetic datasets. This strategy effectively balances multi-domain exposure and enhances the overall depth estimation accuracy and generalization capability in mixed and urban scenes.
\item \textbf{Contribution 3}: We establish a quantitative theoretical verification system for the physical root causes of EPI assumption failure, shifting the research focus from network parameter tuning to physical hypothesis validation. By systematically calculating metrics such as the cone score and angular gradient, we quantitatively assess the degree of angular cone constraint violation across different material domains. This empirical pipeline demonstrates that specular and scattering surfaces inherently disrupt EPI line structures, theoretically falsifying the feasibility of overcoming physical assumption violations solely through architectural or loss function modifications, and thereby guiding future research toward non-EPI paradigms.
\end{itemize}

Extensive experiments demonstrate that our proposed method achieves an overall Mean Absolute Error (MAE) of 0.133 and a mixed/urban scene MAE of 0.081, significantly outperforming baseline EPI-based methods in complex environments. Concurrently, the quantitative analysis confirms that while EPI-based architectures remain highly effective for Lambertian and mixed scenes, they encounter an inherent physical limitation in purely non-Lambertian domains due to the breakdown of the angular cone constraint.

\section{Related Work}

\subsection{Traditional Light Field Depth Estimation Methods}

Traditional light field depth estimation predominantly relies on the geometric properties of Epipolar Plane Images (EPIs). Wanner et al. \cite{he2016} formulated depth inference as an orientation estimation problem by computing the structure tensor of EPIs. The dominant eigenvector of the structure tensor indicates the slope of the linear structures, which is inversely proportional to the scene depth. Subsequent extensions integrated variational frameworks to enforce spatial smoothness and global consistency. While these geometric approaches yield accurate depth maps for matte surfaces, their efficacy is intrinsically bound to the integrity of EPI line structures, assuming constant radiance across viewpoints.

To address occlusions and textureless regions where EPI lines become ambiguous, researchers have integrated defocus and correspondence cues. Tao et al. \cite{vaswani2017} proposed a unified framework that combines the depth from defocus (DfD) and depth from correspondence (DfC) by evaluating the variance of focal stacks and the photo-consistency across angular views, respectively. By fusing the cost volumes derived from both cues, this method enhances depth estimation in occluded areas. Jeon et al. advanced this paradigm by incorporating multi-view stereo (MVS) techniques with sub-pixel phase-based shift estimation to improve spatial resolution. Although fusing multiple cues improves robustness against occlusions, these methods still operate under the implicit assumption that angular photo-consistency holds for all surface types, limiting their applicability to complex reflectance environments.

Recognizing that surface reflectance properties alter angular light transport, several studies have explored frequency-domain analysis to jointly estimate depth and material properties. By applying 2D Discrete Fourier Transform (2D-DFT) or analyzing the angular frequency spectrum of EPIs, methods such as those proposed by Zhang et al. attempt to separate diffuse and specular components. The amplitude and phase distributions in the frequency domain are utilized to classify materials and refine depth maps, drawing analogies to magnetic resonance imaging (MRI) signal separation. These frequency-based techniques demonstrate improved performance in synthetic or high-angular-resolution datasets by isolating non-Lambertian reflections, providing an initial pathway to incorporate material awareness into depth estimation.

Despite their structural and algorithmic advancements, these traditional methods share significant limitations when applied to complex real-world environments. First, these approaches heavily rely on the Lambertian assumption. When surfaces exhibit specular reflection or scattering (non-Lambertian properties), EPI lines fracture or manifest as X-shaped bimodal distributions, causing traditional slope calculations to fail completely. Second, conventional frequency-domain analysis methods cannot reliably extract features under low angular resolutions (e.g., $9 \times 9$ discrete sampling), leading to a sharp degradation in material classification and depth estimation performance. Third, existing techniques lack a physical-level decoupling capability for depth and material information. Consequently, they cannot effectively distinguish Lambertian from non-Lambertian regions at the pixel level, making it difficult to provide reliable prior masks. These deficiencies indicate that the failure in non-Lambertian domains stems from the breakdown of underlying physical assumptions rather than improper network parameter or loss function design. This physical root cause necessitates a unified architecture capable of physical-level angular signal decomposition and dual-mask routing to explicitly handle mixed reflectance properties.

\subsection{Deep Learning-based General Light Field Depth Estimation}

Shin et al. \cite{vaswani2017} proposed EPINet, which formulates depth estimation as a regression task by feeding multi-directional EPIs into a convolutional neural network (CNN). The network extracts the slope of linear structures within EPIs to infer disparity. Subsequent variants enhanced this paradigm by incorporating multi-scale EPI extraction or pseudo-4D input structures to enrich feature representations. These approaches yield accurate depth maps in Lambertian regions, as their feature extractors are specifically optimized for the linear geometry of EPIs. However, confining feature extraction to fixed convolutional kernels restricts their adaptability when EPI structures are distorted by non-Lambertian reflections.

To exploit the intrinsic 4D structure of light fields more comprehensively, architectures such as LFNet \cite{he2016} and DistgNet introduced specialized spatial-angular convolutions and distance-guided feature aggregation modules. LFNet employs alternating spatial and angular convolutions to separately capture spatial textures and view correlations, whereas DistgNet disentangles spatial and angular features to reduce computational redundancy and enhance multi-view consistency. Although these designs improve feature aggregation, they implicitly learn photo-consistency across views without distinguishing the physical differences between Lambertian diffuse reflection and specular highlights at the feature level.

More recently, attention mechanisms and Transformer architectures have been integrated to capture long-range dependencies and global context within light fields. Several studies utilize self-attention along the angular dimension to dynamically weight and aggregate multi-view features \cite{ronneberger2015}, while others combine multi-task learning to jointly optimize depth and surface normal estimation \cite{he2016}. These advanced architectures improve global consistency in complex scenes. Nevertheless, the allocation of attention weights remains heavily dependent on the quality of input features. When the input contains noise that violates physical assumptions due to non-Lambertian reflections, the attention mechanism tends to erroneously amplify these anomalous responses.

While the aforementioned deep learning methods have substantially advanced general light field depth estimation, they share several intrinsic limitations. First, these networks globally enforce the EPI linearity assumption or photo-consistency constraints, without considering the physical root cause of EPI assumption failure in non-Lambertian regions (i.e., the breakdown of the angular cone geometric constraint). This oversight causes errors in non-Lambertian areas to be amplified by the network. Second, existing architectures lack explicit material awareness or dual-branch routing mechanisms, forcing the network to blindly learn noise that violates physical assumptions, thereby failing to adaptively process complex reflective surfaces in mixed scenes. Third, when confronted with the extreme scarcity of non-Lambertian training data, which typically exhibits a long-tail distribution, conventional training strategies are highly susceptible to severe overfitting. They lack an effective domain-balanced sampling mechanism to maintain exposure balance across multiple domains. Addressing these physical and data-level deficiencies forms the core motivation of the proposed unified dual-mask physical model.

\subsection{Deep Learning-based Non-Lambertian Handling and Physics-Aware Methods}

To mitigate the adverse effects of non-Lambertian reflections, early deep learning approaches focused on explicitly detecting and masking specular regions. Methods leveraging angular variance and high-frequency Epipolar Plane Image (EPI) analysis attempt to identify non-Lambertian pixels by calculating the intensity variance $\sigma^2 = \frac{1}{|\mathcal{V}|} \sum_{\mathbf{v} \in \mathcal{V}} (I(\mathbf{x}, \mathbf{v}) - \bar{I})^2$ across angular views $\mathcal{V}$. These networks typically employ a pre-processing module to generate a specularity mask, which is then used to down-weight or inpaint distorted EPI regions before depth regression. While this masking strategy improves depth accuracy on glossy surfaces, it heavily relies on high angular resolution to reliably compute angular statistics. In sparsely sampled light fields, the discrete angular gradients become ambiguous, causing high-frequency analysis to fail in distinguishing between texture edges and specular highlights.

Moving beyond simple masking, subsequent research explored joint material decomposition and depth estimation to inherently address reflectance variations. Physics-aware architectures integrate intrinsic image decomposition into the depth regression pipeline, simultaneously predicting diffuse albedo, specular reflectance, and depth maps. By enforcing physical consistency constraints, such as the dichromatic reflection model $I(\mathbf{x}, \mathbf{v}) = I_d(\mathbf{x}) + I_s(\mathbf{x}, \mathbf{v})$, these networks learn to separate the view-invariant diffuse component $I_d$ from the view-dependent specular component $I_s$. The extracted diffuse features are then fed into the depth estimator, theoretically restoring the Lambertian assumption. This joint optimization often suffers from ill-posedness, especially when the diffuse and specular components share similar spatial frequencies. These models also lack a rigorous quantitative analysis of why the underlying EPI assumptions fail, treating the non-Lambertian issue merely as a feature entanglement problem rather than a core breakdown of the angular cone constraint.

To adaptively process heterogeneous surface properties without explicit physical decomposition, recent advancements have introduced dual-branch or attention-based routing mechanisms. Architectures employing implicit masking or soft-attention modules process spatial textures and angular geometries in parallel. A learnable gating mechanism dynamically adjusts the contribution of EPI-based geometric features versus spatial context features based on the local reflectance properties. In Lambertian regions, the network relies on EPI slopes, whereas in non-Lambertian regions, it shifts reliance to spatial context or learned priors. Although this implicit routing enhances robustness across mixed scenes, the routing weights are learned entirely in a data-driven manner without explicit physical or geometric priors. Consequently, the network struggles to generalize to unseen material distributions, and the absence of domain-balanced sampling exacerbates the performance degradation in long-tail non-Lambertian scenarios.

Despite these advancements, existing deep learning-based non-Lambertian handling and physics-aware methods exhibit distinct limitations. First, current material classification and specularity detection techniques predominantly depend on high-frequency domain analysis or require dense angular sampling, lacking a robust physical decomposition model based on geometric angular features---such as angular gradients and angular correlations---under low angular resolution conditions. Second, prior works fail to quantitatively demonstrate the root cause of EPI failure from a physical perspective, specifically the violation of the angular cone constraint. This absence of a theoretical closed-loop for hypothesis verification and dead-end identification leads to blind network architecture design and inefficient trial-and-error processes. Finally, existing dual-branch or masking approaches rely on implicit learning, lacking an explicit dual-mask routing architecture guided by high-precision physical and geometric prior masks. They also do not effectively integrate domain-balanced sampling to overcome the generalization bottleneck caused by the scarcity of non-Lambertian training data. Addressing these gaps, the proposed unified dual-mask physical model introduces a three-layer angular signal physical decomposition to robustly decouple depth and material, a quantitative theoretical verification system for EPI hypothesis failure, and an explicit dual-mask routing architecture with domain-balanced sampling to achieve reliable depth estimation in complex real-world environments.

While existing methods mitigate non-Lambertian effects by explicitly masking specular regions via angular variance and high-frequency Epipolar Plane Image (EPI) analysis, they heavily rely on dense angular sampling and blindly assume the universal validity of EPI linearity, leading to severe performance degradation under sparse views and complex physical reflections. Inspired by these observations, we propose a unified dual-mask physical approach that addresses the aforementioned limitations by quantitatively revealing the physical root causes of EPI failures, decoupling depth and material via a three-layer angular signal decomposition, and leveraging dual-mask routing with domain-balanced sampling to ensure robust depth estimation across complex non-Lambertian and long-tail scenarios. Building on these foundational insights, we now formalize the proposed approach into a comprehensive computational framework, which is detailed in the subsequent section.

\section{Methodology}

In this section, we present the overall architecture of the proposed Unified Dual-Mask Physical Model for non-Lambertian light field depth estimation. Different from previous works that implicitly assume a uniform Lambertian reflectance across the entire scene, our framework explicitly models the physical heterogeneity of real-world surfaces. By integrating a physical signal decomposition mechanism with a dual-branch routing strategy, the proposed architecture effectively addresses the failure of Epipolar Plane Image (EPI) linear assumptions on specular and scattering materials. This design not only enhances depth accuracy in complex multi-view imaging and light field video systems but also provides a theoretically grounded solution for domain-specific feature extraction. As illustrated in Fig. 1, the network comprises a Three-Layer Angular Signal Decomposition module, a dual-branch feature extraction pipeline (Lambertian EPI Branch and Non-Lambertian Geometric Branch), a Per-pixel Geometric Mask Generator, and a Mask-Weighted Fusion \& Depth Regression module.

The input to the network is a light field image represented as a 4D tensor, from which we extract 2D angular sampling patches. Specifically, for each spatial pixel $(x, y)$, the input is formulated as an angular patch $I(u,v) \in \mathbb{R}^{U \times V \times C}$, where $u$ and $v$ denote the angular coordinates, and $U \times V$ represents the angular resolution (e.g., $9 \times 9$ in our implementation). $C$ indicates the number of color channels. Prior to feeding into the main network, the input patches undergo standard normalization to stabilize the training process. Concurrently, we compute the angular gradient $\nabla_{\theta} I$ across the angular dimensions to capture the local geometric variations, which serves as a primary indicator for distinguishing Lambertian from non-Lambertian regions. These preprocessed angular patches and their corresponding gradient maps are then routed in parallel to the decomposition and mask generation modules.

The core data flow begins with the \textbf{Three-Layer Angular Signal Decomposition} module, which is motivated by the need to decouple the mixed physical attributes inherent in light field signals. Unlike conventional methods that process the raw angular signal holistically, our module explicitly decomposes $I(u,v)$ into three distinct physical components:
\begin{equation}
I(u,v) = DC + a \cdot u + b \cdot v + \varepsilon(u,v)
\end{equation}
where $DC$ represents the ambient illumination (direct current component), $a$ and $b$ are the linear coefficients corresponding to the EPI slopes induced by parallax and depth, and $\varepsilon(u,v)$ denotes the residual component capturing material reflectance properties. This decomposition yields the illumination and depth-related features for the \textbf{Lambertian EPI Branch}, while the residual $\varepsilon(u,v)$ is directed to the \textbf{Non-Lambertian Geometric Branch}. In the Lambertian EPI Branch, the objective is to extract smooth depth features for matte surfaces that satisfy the angular cone constraint. Utilizing the linear coefficients $a$ and $b$, this branch employs a Plane Regular Sampling Operator (PRSO) or a displacement field to aggregate features along the EPI lines, outputting the Lambertian depth feature map $F_L$. Conversely, the Non-Lambertian Geometric Branch is designed to handle surfaces that violate the EPI linearity assumption. Instead of relying on frequency-domain analysis which degrades at low angular resolutions, this branch leverages a Geometric Extractor to process the angular gradient $\nabla_{\theta} I$ and the residual $\varepsilon(u,v)$. It captures geometric angular features, such as bimodal distributions and X-shaped patterns, producing the non-Lambertian depth feature map $F_{NL}$. Both $F_L$ and $F_{NL}$ are subsequently projected to a unified channel dimension ($\text{FUSE\_DIM} = 96$) to ensure balanced representation. Simultaneously, the \textbf{Per-pixel Geometric Mask Generator} processes the angular gradient and residual features to produce domain-specific routing masks. Motivated by the high noise levels associated with scene-level binary cross-entropy supervision, this module utilizes a Multi-Layer Perceptron (MLP) to generate pixel-wise pseudo-labels. The non-Lambertian mask $M_{NL}$ and Lambertian mask $M_L$ are computed as:
\begin{multline}

M_{NL} = \sigma(\text{MLP}(\nabla_{\theta} I, \varepsilon)), \quad M_L \\
= 1 - M_{NL}
\end{multline}
where $\sigma$ denotes the sigmoid activation function. This mechanism provides highly discriminative, pixel-level material classification, effectively guiding the subsequent feature fusion.

The final stage of the architecture is the \textbf{Mask-Weighted Fusion \& Depth Regression} module, which aims to integrate the domain-specific features while eliminating cross-domain interference. Different from previous works that simply concatenate or add multi-branch features, our module explicitly applies the generated dual masks to route the features via element-wise multiplication (Hadamard product). The fused feature representation is formulated as:
\begin{equation}
F_{fused} = M_L \odot F_L + M_{NL} \odot F_{NL}
\end{equation}
where $\odot$ denotes the Hadamard product. This mask-weighted aggregation ensures that the depth estimation for each pixel is predominantly driven by the branch corresponding to its physical material property. The fused features $F_{fused}$ are then passed through a regression head to predict the final high-precision depth map $D_{final}$, alongside a domain classification prediction. The output $D_{final}$ retains the spatial resolution of the input central view, providing a dense and accurate depth estimation that is robust to the complex reflectance properties encountered in real-world light field video and multi-view systems.


\begin{figure*}[!t]
\centering
\resizebox{\textwidth}{!}{
\begin{tikzpicture}[
  module/.style={
    draw, thick, rounded corners=4pt, fill=blue!15,
    minimum height=1.8cm, text width=2.2cm, align=center,
    font=\small\bfseries, outer sep=3pt, inner sep=4pt
  },
  data/.style={
    draw, thick, shape=parallelogram, shape border rotate=0,
    minimum height=1.8cm, text width=1.8cm, align=center,
    font=\small\bfseries, outer sep=3pt, inner sep=4pt, fill=gray!15
  },
  innov1/.style={fill=orange!30},
  innov2/.style={fill=red!20},
  arr/.style={-{Stealth[length=2.5mm]}, thick, color=black!80},
  formula/.style={font=\scriptsize\itshape, text=black!70}
]

  % --- Nodes ---
  \node (input) [data] {Light Field Images\\{\scriptsize (9$\times$9 Angular)}};
  
  \node (decomp) [module, innov1, right=1.2cm of input] {Three-Layer Angular\\Signal Decomposition\\{\scriptsize $I=DC+au+bv+\varepsilon$}};
  
  \node (lambert) [module, innov2, right=1.2cm of decomp] {Lambertian\\EPI Branch\\{\scriptsize PRSO / Disp. Field}};
  
  \node (fusion) [module, innov1, right=1.2cm of lambert, yshift=-1.5cm] {Mask-Weighted Fusion\\\\& Depth Regression\\{\scriptsize $D=M_L F_L + M_{NL} F_{NL}$}};
  
  \node (output) [data, right=1.2cm of fusion] {Final Depth Map\\\\& Domain Mask};

  \node (mask_gen) [module, innov1, below=2.8cm of decomp] {Per-pixel Geometric\\Mask Generator\\{\scriptsize $\sigma(\text{MLP}(\nabla_\theta I))$}};
  
  \node (non_lambert) [module, innov1, right=1.2cm of mask_gen] {Non-Lambertian\\Geometric Branch\\{\scriptsize Angular Gradient}};

  % --- Connections ---
  % Input splits
  \coordinate (in_split) at ([xshift=0.6cm]input.east);
  \draw[arr] (input.east) -- (in_split);
  \draw[arr] (in_split) |- (decomp.west);
  \draw[arr] (in_split) |- (mask_gen.west);

  % Decomp splits
  \coordinate (dec_split) at ([xshift=0.6cm]decomp.east);
  \draw[arr] (decomp.east) -- (dec_split);
  \draw[arr] (dec_split) |- node[pos=0.25, above, font=\scriptsize] {DC, $a, b$} (lambert.west);
  \draw[arr] (dec_split) |- node[pos=0.25, above, font=\scriptsize] {Residual $\varepsilon$} (non_lambert.west);

  % Branches to Fusion
  \draw[arr] (lambert.east) -- ++(0.5,0) |- node[pos=0.25, above, font=\scriptsize] {$F_L$} (fusion.north);
  \draw[arr] (non_lambert.east) -- ++(0.5,0) |- node[pos=0.25, below, font=\scriptsize] {$F_{NL}$} ([yshift=3mm]fusion.south);

  % Mask Gen to Fusion
  \draw[arr] (mask_gen.south) -- ++(0,-0.6) -| node[pos=0.25, right, font=\scriptsize] {$M_L, M_{NL}$} ([yshift=-3mm]fusion.south);

  % Fusion to Output
  \draw[arr] (fusion.east) -- (output.west);

  % --- Legend ---
  \node[font=\scriptsize, anchor=west] (leg1) at ([yshift=-0.8cm]current bounding box.south west) {\tikz\draw[fill=orange!30, thick] (0,0) rectangle (0.25,0.25); Physical Decomposition \\& Dual-Mask};
  \node[font=\scriptsize, anchor=west, right=0.3cm of leg1] (leg2) {\tikz\draw[fill=red!20, thick] (0,0) rectangle (0.25,0.25); EPI Failure Analysis};
  \node[font=\scriptsize, anchor=west, right=0.3cm of leg2] (leg3) {\tikz\draw[fill=blue!15, thick] (0,0) rectangle (0.25,0.25); Standard Module};
  \node[font=\scriptsize, anchor=west, right=0.3cm of leg3] (leg4) {\tikz\draw[fill=gray!15, thick] (0,0) rectangle (0.25,0.25); Data};

\end{tikzpicture}
}
\caption{Overall architecture of the proposed method.}
\label{fig:architecture}
\end{figure*}

\subsection{Three-Layer Angular Signal Decomposition}

Following the extraction of the 2D angular sampling patches $I(u,v) \in \mathbb{R}^{U \times V \times C}$, the network processes these inputs through the Three-Layer Angular Signal Decomposition module. Conventional light field depth estimation methods typically assume a uniform Lambertian reflectance across the scene, relying heavily on the linear structure of Epipolar Plane Images (EPIs). In real-world multi-view imaging and light field video systems, non-Lambertian surfaces, such as specular and scattering materials, violate this linear assumption, leading to severe depth estimation artifacts. Existing deep learning approaches often implicitly learn these features within black-box architectures, lacking physical interpretability and struggling to generalize across diverse material types. To address this limitation, we design the Three-Layer Angular Signal Decomposition module to explicitly decouple the physical attributes of the light field signal. The primary objective of this module is to separate the angular signal into illumination, geometry, and material components, providing physically disentangled feature representations for the subsequent dual-branch processing.

The architecture of this module is grounded in the physical rendering equation and surface reflectance models. For a specific spatial pixel and a given channel $c$, the observed angular signal $I_c(u,v)$ is modeled as a superposition of ambient illumination, linear displacement induced by parallax, and high-frequency residuals caused by non-Lambertian reflections. This physical decomposition is formulated as:
\begin{equation}
I_c(u,v) = \mu_c + a_c \cdot u + b_c \cdot v + \varepsilon_c(u,v)
\end{equation}
where $\mu_c$ represents the direct current (DC) component, which captures the ambient illumination and base albedo. The terms $a_c$ and $b_c$ are the linear coefficients corresponding to the horizontal and vertical EPI slopes, respectively, directly encoding the scene parallax and depth information. The term $\varepsilon_c(u,v)$ denotes the residual component, which isolates the non-Lambertian material characteristics, such as specular highlights and anisotropic reflections, that deviate from the linear EPI assumption.

To robustly estimate these parameters from discrete angular samples, we employ a least-squares optimization to fit the planar model. The objective function for channel $c$ is defined as:
\begin{equation}
\begin{aligned}
\mathcal{L}_{dec}(\mu_c, a_c, b_c) &= \sum_{u=1}^{U} \sum_{v=1}^{V} \left( I_c(u,v) \right. \\
&\quad \left. - (\mu_c + a_c u + b_c v) \right)^2
\end{aligned}
\end{equation}
By setting the partial derivatives with respect to $\mu_c$, $a_c$, and $b_c$ to zero, we obtain a closed-form solution. This optimization is efficiently expressed in matrix notation as:
\begin{equation}
\mathbf{Y}_c = \mathbf{X} \boldsymbol{\beta}_c + \boldsymbol{\varepsilon}_c
\end{equation}
where $\mathbf{Y}_c \in \mathbb{R}^{UV \times 1}$ is the flattened angular signal vector, $\mathbf{X} \in \mathbb{R}^{UV \times 3}$ is the design matrix containing the coordinates $[1, u, v]$ for each angular sample, and $\boldsymbol{\beta}_c = [\mu_c, a_c, b_c]^T$ is the parameter vector. The optimal parameter estimate is computed as:
\begin{multline}

\hat{\boldsymbol{\beta}}_c \\
= (\mathbf{X}^T \mathbf{X})^{-1} \mathbf{X}^T \mathbf{Y}_c
\end{multline}
Since the design matrix $\mathbf{X}$ depends exclusively on the angular coordinates $(u,v)$, the projection matrix $(\mathbf{X}^T \mathbf{X})^{-1} \mathbf{X}^T$ can be precomputed. This property renders the decomposition process computationally efficient, making it highly suitable for high-resolution light field video streams.

Following the estimation of the DC and linear components, the residual component is calculated by subtracting the fitted planar model from the original angular signal:
\begin{equation}
\hat{\varepsilon}_c(u,v) = I_c(u,v) - (\hat{\mu}_c + \hat{a}_c u + \hat{b}_c v)
\end{equation}
The module outputs the DC component $\hat{\boldsymbol{\mu}} \in \mathbb{R}^{C}$, the linear coefficients $\hat{\mathbf{a}}, \hat{\mathbf{b}} \in \mathbb{R}^{C}$, and the residual feature map $\hat{\boldsymbol{\varepsilon}} \in \mathbb{R}^{U \times V \times C}$. These outputs maintain the channel dimension $C$ while explicitly separating the physical properties across different tensor structures.

The disentangled outputs establish a direct connection to the subsequent dual-branch architecture. The DC component $\hat{\boldsymbol{\mu}}$ and linear coefficients $\hat{\mathbf{a}}, \hat{\mathbf{b}}$ are routed to the Lambertian EPI Branch, which utilizes the continuous depth prior to extract smooth depth features for regions adhering to the angular cone constraint. Concurrently, the residual feature map $\hat{\boldsymbol{\varepsilon}}$, along with the original input, is directed to the Non-Lambertian Geometric Branch to extract geometric angular features for complex reflective surfaces. The residual features also serve as an input to the Per-pixel Geometric Mask Generator. Different from previous works that process the light field signal as an entangled entity, our module explicitly introduces a physical signal decomposition mechanism. This explicit decoupling reduces the learning burden on the subsequent networks, mitigates the interference of non-Lambertian regions on EPI slope estimation, and enhances the generalization capability of the overall system across diverse and complex material types.

\subsection{Lambertian EPI Branch}

Following the physical decoupling in the Three-Layer Angular Signal Decomposition module, the light field signal is separated into illumination, geometry, and material components. For Lambertian surfaces, the angular signal strictly adheres to the photo-consistency constraint, manifesting as linear structures in Epipolar Plane Images (EPIs). The primary objective of the Lambertian EPI Branch is to extract robust and smooth depth features for these regions by leveraging the continuous depth prior derived from the linear coefficients. Conventional light field depth estimation methods typically rely on discrete disparity search or frequency-domain analysis along EPI lines. These approaches are inherently susceptible to quantization errors and struggle in textureless or occluded regions. Unlike prior works that implicitly learn EPI structures through black-box convolutions, our module explicitly utilizes the physically disentangled linear coefficients to guide a continuous resampling process, thereby ensuring geometric consistency and sub-pixel accuracy.

The architecture of this branch takes the linear coefficients $a, b \in \mathbb{R}^{H \times W}$ and the extracted EPI features $F_{in} \in \mathbb{R}^{H \times W \times U \times V \times C_{in}}$ as inputs, where $H \times W$ denotes the spatial resolution, $U \times V$ represents the angular resolution, and $C_{in}$ is the feature channel dimension. The linear coefficients $a(\mathbf{p})$ and $b(\mathbf{p})$ at a spatial location $\mathbf{p} = (x,y)$ correspond to the horizontal and vertical EPI slopes, respectively. We first formulate the continuous disparity vector field $\mathbf{d}(\mathbf{p})$:
\begin{multline}

\mathbf{d}(\mathbf{p}) \\
= \begin{bmatrix} a(\mathbf{p}) \\ b(\mathbf{p}) \end{bmatrix}
\end{multline}
This vector field serves as the continuous depth prior, dictating the geometric displacement of the EPI lines across different views.

To exploit the linear EPI structure without the limitations of discrete disparity hypotheses, we employ a Planar Regularized Sampling Operator (PRSO) combined with a differentiable displacement field. For a given spatial coordinate $\mathbf{p}$ and an angular coordinate $\mathbf{u} = (u,v) \in \Omega$, where $\Omega$ is the set of angular samples, the theoretical epipolar projection $\mathbf{p}'(\mathbf{p}, \mathbf{u})$ is calculated as:
\begin{equation}
\begin{aligned}
\mathbf{p}'(\mathbf{p}, \mathbf{u}) &= \mathbf{p} + \mathbf{d}(\mathbf{p}) \odot \mathbf{u} \\
&= \begin{bmatrix} x + u \cdot a(\mathbf{p}) \\ y + v \cdot b(\mathbf{p}) \end{bmatrix}
\end{aligned}
\end{equation}
where $\odot$ denotes the element-wise multiplication. The PRSO then samples the input EPI features at these continuous sub-pixel coordinates. Mathematically, the resampled feature $\tilde{F}(\mathbf{p}, \mathbf{u})$ is defined as:
\begin{equation}
\begin{aligned}
\tilde{F}(\mathbf{p}, \mathbf{u}) &= \text{PRSO}(F_{in}(\cdot, \mathbf{u}), \\
&\quad \mathbf{p}'(\mathbf{p}, \mathbf{u})) \\
&= \sum_{\mathbf{q} \in \mathcal{N}(\mathbf{p}')} F_{in}(\mathbf{q}, \mathbf{u}) \\
&\quad \cdot K(\mathbf{q} - \mathbf{p}'(\mathbf{p}, \mathbf{u}))
\end{aligned}
\end{equation}
where $\mathcal{N}(\mathbf{p}')$ represents the spatial neighborhood of the projected coordinate, and $K(\cdot)$ is a differentiable interpolation kernel (e.g., bilinear or bicubic kernel) regularized to preserve high-frequency structural details. This continuous resampling effectively aligns the angular features along the EPI lines, compensating for the parallax induced by depth.

Following the geometric alignment, the branch aggregates the resampled features across the angular dimension to enhance the photo-consistency signal and suppress residual noise. The aggregated feature map $F_{agg} \in \mathbb{R}^{H \times W \times C_{in}}$ is computed by averaging the aligned features over all angular views:
\begin{multline}

F_{agg}(\mathbf{p}) \\
= \frac{1}{|\Omega|} \sum_{\mathbf{u} \in \Omega} \tilde{F}(\mathbf{p}, \mathbf{u})
\end{multline}
where $|\Omega| = U \times V$ is the total number of angular samples. This aggregation operation projects the 4D light field features into a 2D spatial representation, concentrating the depth information for Lambertian regions.

To ensure dimensional compatibility and balance the feature distribution for the subsequent fusion stage, the aggregated features are passed through a FUSE\_DIM projection layer. This layer consists of a $1 \times 1$ convolution followed by a normalization and activation function, mapping the features to a unified channel space $C_{out}$. The final Lambertian depth feature map $F_L \in \mathbb{R}^{H \times W \times C_{out}}$ is formulated as:
\begin{equation}
\begin{aligned}
F_L &= \text{FUSE\_DIM}(F_{agg}) \\
&= \sigma(\mathbf{W}_{fuse} * F_{agg} + \mathbf{b}_{fuse})
\end{aligned}
\end{equation}
where $\mathbf{W}_{fuse}$ and $\mathbf{b}_{fuse}$ are the learnable weights and biases of the projection layer, $*$ denotes the convolution operation, and $\sigma$ is the activation function.

The proposed Lambertian EPI Branch explicitly integrates physical EPI constraints with continuous geometric priors, distinguishing it from conventional cost-volume or spectrum-based methods. By utilizing the continuous displacement field derived from the physical decomposition, the module avoids the quantization artifacts inherent in discrete disparity search and provides sub-pixel accurate depth features. The resulting feature map $F_L$ encapsulates highly reliable and smooth depth representations for Lambertian surfaces, which is subsequently routed to the Mask-Weighted Fusion module to be integrated with the non-Lambertian geometric features.

\subsection{Non-Lambertian Geometric Branch}

While the Lambertian EPI Branch effectively processes regions adhering to the photo-consistency constraint, real-world scenes extensively contain non-Lambertian surfaces, such as specular highlights, reflections, and transparent materials. These surfaces violate the angular cone constraint, causing severe distortions in Epipolar Plane Images (EPIs), such as bimodal distributions or X-shaped patterns. Conventional light field depth estimation methods typically rely on frequency-domain analysis or discrete disparity search along EPI lines. These approaches struggle to model the complex angular variations of non-Lambertian regions, often yielding severe depth artifacts. Different from previous works that implicitly learn non-Lambertian features through black-box convolutions or apply frequency-domain filters, our Non-Lambertian Geometric Branch explicitly extracts geometric angular features to robustly estimate depth in these challenging regions.

The architecture of this branch takes the original light field angular sampling images $I \in \mathbb{R}^{H \times W \times U \times V \times C_{in}}$ and the residual component $\varepsilon \in \mathbb{R}^{H \times W \times U \times V \times C_{res}}$ decomposed by the Three-Layer Angular Signal Decomposition module as inputs. Here, $H \times W$ denotes the spatial resolution, $U \times V$ represents the angular resolution, and $C_{in}, C_{res}$ are the respective channel dimensions. The residual component $\varepsilon(u,v)$ encapsulates the material and reflectance properties decoupled from illumination and geometry. To capture the distinct structural variations caused by non-Lambertian reflectance, we first compute the angular gradient $\nabla_{\theta} I$ across the angular dimensions $u$ and $v$. For a specific spatial location $(x,y)$, the angular gradient is formulated as:
\begin{equation}
\begin{aligned}
\nabla_{\theta} I(x,y,u,v) &= \left[ \frac{\partial I(x,y,u,v)}{\partial u}, \right. \\
&\quad \left. \frac{\partial I(x,y,u,v)}{\partial v} \right]
\end{aligned}
\end{equation}
where the partial derivatives are approximated using central finite differences along the angular grid. This gradient highlights the abrupt intensity changes characteristic of specularities and occlusions.

To model the specific EPI patterns of non-Lambertian surfaces, we design a Geometric Extractor that integrates the angular gradient with the physical residual component. Unlike Lambertian surfaces that exhibit linear EPI structures, non-Lambertian surfaces generate X-shaped patterns at occlusion boundaries or bimodal distributions under specular reflections. We capture these patterns by computing the local angular structure tensor. Let $\Omega$ be a local angular window centered at $(u,v)$. The angular covariance matrix $\mathbf{J}(x,y,u,v)$ is computed as:
\begin{multline}
\mathbf{J}(x,y,u,v) = \sum_{(u',v') \in \Omega} w(u',v') \\
\times \nabla_{\theta} I(x,y,u',v')^T \nabla_{\theta} I(x,y,u',v')
\end{multline}
where $w(u',v')$ is a Gaussian weighting function. The eigenvalues $\lambda_1, \lambda_2$ of $\mathbf{J}$ indicate the dominant angular variation directions, effectively distinguishing X-shaped crossings (where both eigenvalues are large) from linear structures (where one eigenvalue dominates). We then concatenate the structural tensor features with the residual component $\varepsilon(u,v)$ to form a comprehensive geometric representation:
\begin{equation}
\begin{aligned}
F_{geom} &= \text{Conv}_{geo} \left( \text{Concat}(\lambda_1, \lambda_2, \right. \\
&\quad \left. \mathbf{v}_1, \varepsilon(u,v)) \right)
\end{aligned}
\end{equation}
where $\mathbf{v}_1 \in \mathbb{R}^2$ is the eigenvector corresponding to $\lambda_1$, and $\text{Conv}_{geo}$ denotes a series of 3D convolutional layers that aggregate spatial-angular contextual information and reduce the angular dimensions, yielding $F_{geom} \in \mathbb{R}^{H \times W \times C_{geo}}$. This explicit geometric formulation allows the network to directly target non-Lambertian anomalies rather than treating them as noise.

To ensure compatibility with the Lambertian EPI Branch and eliminate domain shift between the two feature spaces, the extracted geometric features $F_{geom}$ are passed through a Feature Unified Spatial-Epipolar Dimensional (FUSE\_DIM) projection module. The FUSE\_DIM module aligns the channel dimensions and normalizes the feature distributions, yielding the final non-Lambertian depth feature map:
\begin{equation}
F_{NL} = \text{FUSE\_DIM}(F_{geom})
\end{equation}
where $F_{NL} \in \mathbb{R}^{H \times W \times C_{out}}$ represents the projected feature map with the output channel dimension $C_{out}$, matching the dimension of the Lambertian feature map $F_L$.

The Non-Lambertian Geometric Branch operates in parallel with the Lambertian EPI Branch. It receives the decoupled residual component from the Three-Layer Angular Signal Decomposition module and outputs $F_{NL}$, which is subsequently routed to the Mask-Weighted Fusion and Depth Regression module. By explicitly leveraging angular gradients and physical residuals to model bimodal and X-shaped EPI patterns, this branch provides reliable depth cues for regions where traditional photo-consistency assumptions fail, thereby enhancing the overall robustness of the unified dual-mask physical model.

\subsection{Per-pixel Geometric Mask Generator}

Dual-branch architectures for light field depth estimation require accurate domain masks to separate Lambertian and non-Lambertian regions. Conventional methods typically employ scene-level binary cross-entropy (BCE) supervision or learn domain classification implicitly through black-box classifiers. Scene-level labels contain high noise and fail to capture pixel-wise reflectance variations, leading to boundary artifacts and domain confusion during feature fusion. To address this limitation, we introduce the Per-pixel Geometric Mask Generator, which explicitly derives pixel-wise masks from physical geometric properties to guide the dual-branch feature integration.

Different from previous works that rely on coarse scene-level BCE supervision or implicit feature clustering, our module leverages geometric features with high statistical discriminability (Cohen's d = 0.983) to generate pixel-wise pseudo-labels and masks. This provides a physically interpretable and pixel-accurate domain classification mechanism, avoiding the gradient conflicts caused by noisy global labels.

The architecture of this module takes the angular gradient feature $G \in \mathbb{R}^{H \times W \times C_g}$ and the residual component $\varepsilon \in \mathbb{R}^{H \times W \times U \times V \times C_{res}}$ as inputs. Here, $H \times W$ denotes the spatial resolution, $U \times V$ represents the angular resolution, and $C_g, C_{res}$ are the respective channel dimensions. The angular gradient feature $G$ captures the local angular variations, while the residual component $\varepsilon$ encapsulates the material-specific deviations decomposed by the Three-Layer Angular Signal Decomposition module.

To align the spatial and channel dimensionalities, the residual component $\varepsilon$ is aggregated along the angular dimensions $(U, V)$ to form a spatial residual feature $R \in \mathbb{R}^{H \times W \times C_r}$. The aggregation operation computes the angular variance to highlight non-Lambertian deviations:
\begin{equation}
\begin{aligned}
R(x,y,c) &= \frac{1}{UV} \sum_{u=1}^{U} \sum_{v=1}^{V} \\
&\quad \times \left( \varepsilon(x,y,u,v,c) - \bar{\varepsilon}(x,y,c) \right)^2
\end{aligned}
\end{equation}
where $(x,y)$ denotes the spatial coordinates, $c \in \{1, \dots, C_{res}\}$ is the channel index, and $\bar{\varepsilon}(x,y,c)$ is the angular mean of the residual component defined as:
\begin{equation}
\begin{aligned}
\bar{\varepsilon}(x,y,c) &= \frac{1}{UV} \sum_{u=1}^{U} \sum_{v=1}^{V} \\
&\quad \varepsilon(x,y,u,v,c)
\end{aligned}
\end{equation}

The aggregated residual feature $R$ and the angular gradient feature $G$ are concatenated along the channel dimension to form the unified geometric representation $F_{mask} \in \mathbb{R}^{H \times W \times (C_g + C_r)}$:
\begin{equation}
F_{mask} = \text{Concat}(G, R)
\end{equation}

This representation is then processed by a Multi-Layer Perceptron (MLP) to predict the non-Lambertian probability map. The non-Lambertian mask $M_{NL} \in \mathbb{R}^{H \times W \times 1}$ is computed as:
\begin{equation}
M_{NL} = \sigma(\text{MLP}(F_{mask}))
\end{equation}
where $\sigma(\cdot)$ denotes the Sigmoid activation function, which normalizes the output to the range $[0, 1]$. The MLP consists of three $1 \times 1$ convolutional layers interleaved with ReLU activations, mapping the concatenated features to a single-channel probability map.

Correspondingly, the Lambertian mask $M_L \in \mathbb{R}^{H \times W \times 1}$ is derived by taking the complement of $M_{NL}$:
\begin{equation}
M_L = 1 - M_{NL}
\end{equation}
where $1$ represents a tensor of ones with the same spatial dimensions as $M_{NL}$.

The generated masks $M_L$ and $M_{NL}$ are subsequently fed into the Mask-Weighted Fusion \& Depth Regression module. In this downstream module, the masks perform pixel-wise weighting on the depth features from the Lambertian EPI Branch ($F_L$) and the Non-Lambertian Geometric Branch ($F_{NL}$). By utilizing physically grounded geometric features rather than noisy scene-level labels, the Per-pixel Geometric Mask Generator yields highly accurate domain separation. This explicit masking strategy suppresses cross-domain interference, ensuring that the final depth regression maintains sharp boundaries around specular highlights and transparent objects.

\subsection{Mask-Weighted Fusion \& Depth Regression}

Following the extraction of domain-specific features and the generation of pixel-wise physical masks, the final stage of the unified dual-mask physical model is the Mask-Weighted Fusion and Depth Regression module. In conventional dual-branch architectures for light field depth estimation, feature fusion is typically performed via simple concatenation, element-wise addition, or implicit attention mechanisms. These heuristic fusion strategies often suffer from cross-domain interference, where features from the Lambertian branch inadvertently contaminate the depth prediction in non-Lambertian regions, and vice versa. This interference is particularly detrimental at object boundaries and specular highlights, leading to severe depth blurring and structural artifacts. To address this limitation, we design a mask-weighted fusion mechanism that explicitly isolates and integrates domain-specific representations guided by the physically derived masks, ensuring that each spatial location relies exclusively on its corresponding reflectance-specific features.

The module takes four primary inputs: the Lambertian depth feature map $F_L \in \mathbb{R}^{H \times W \times C_f}$ from the Lambertian EPI branch, the non-Lambertian depth feature map $F_{NL} \in \mathbb{R}^{H \times W \times C_f}$ from the non-Lambertian geometric branch, and the corresponding pixel-wise masks $M_L \in \mathbb{R}^{H \times W \times 1}$ and $M_{NL} \in \mathbb{R}^{H \times W \times 1}$ from the per-pixel geometric mask generator. Here, $H \times W$ denotes the spatial resolution, and $C_f$ represents the unified channel dimension after feature projection. The core operation is the mask-weighted feature integration, which computes the fused feature map $F_{fused} \in \mathbb{R}^{H \times W \times C_f}$ as:
\begin{equation}
F_{fused} = M_L \odot F_L + M_{NL} \odot F_{NL}
\end{equation}
where $\odot$ denotes the Hadamard product. Since $M_L$ and $M_{NL}$ are single-channel spatial maps, they are broadcasted along the channel dimension $C_f$ during the multiplication. For any specific spatial coordinate $(x, y)$, the fusion process is formulated as:
\begin{equation}
F_{fused}(x,y) = M_L(x,y) \cdot F_L(x,y) + M_{NL}(x,y) \cdot F_{NL}(x,y)
\end{equation}
Given the complementary nature of the masks generated in the preceding module, where $M_L(x,y) + M_{NL}(x,y) = 1$, this formulation guarantees a convex combination of the dual-branch features at each pixel. Consequently, the feature representation in purely Lambertian regions is entirely dictated by $F_L$, while non-Lambertian regions rely solely on $F_{NL}$. At transition boundaries, the continuous mask values provide a smooth, physically grounded interpolation between the two domains, effectively eliminating the abrupt feature mutations caused by hard thresholding.

Following the feature fusion, the integrated representation $F_{fused}$ is fed into a depth regression head to predict the final depth map. The regression head, denoted as $\mathcal{R}_{depth}$, consists of a series of $1 \times 1$ convolutional layers followed by a ReLU activation, which projects the high-dimensional features into a single-channel depth map $D_{final} \in \mathbb{R}^{H \times W \times 1}$:
\begin{equation}
D_{final} = \mathcal{R}_{depth}(F_{fused})
\end{equation}
To further refine the domain classification and provide auxiliary supervision during training, a parallel mask prediction head $\mathcal{R}_{mask}$ is employed to output the refined domain probability map $P_{mask} \in \mathbb{R}^{H \times W \times 1}$:
\begin{equation}
P_{mask} = \sigma(\mathcal{R}_{mask}(F_{fused}))
\end{equation}
where $\sigma(\cdot)$ is the sigmoid activation function. The overall training objective for this module combines the depth regression loss $\mathcal{L}_{depth}$ (e.g., Smooth L1 loss) and the mask classification loss $\mathcal{L}_{mask}$ (e.g., Binary Cross-Entropy loss using the geometric pseudo-labels), formulated as $\mathcal{L}_{total} = \mathcal{L}_{depth} + \lambda \mathcal{L}_{mask}$, where $\lambda$ is a balancing hyperparameter.

Different from previous works that implicitly fuse dual-branch features via black-box attention modules or simple concatenation, our module explicitly leverages physically interpretable masks to isolate domain-specific representations. This explicit isolation prevents the gradient conflicts and feature contamination inherent in implicit fusion strategies. By strictly confining the Lambertian EPI features to diffuse regions and the non-Lambertian geometric features to specular or complex reflectance regions, the proposed fusion mechanism substantially enhances the structural integrity of the depth map. The expected effect is a significant reduction in boundary artifacts and an improved depth accuracy in challenging non-Lambertian areas, yielding a robust and high-fidelity light field depth estimation system.

\subsection{Training Objective}

\subsubsection{Training Objective and Loss Function}

\paragraph{Overview of the Joint Optimization Framework}
The unified dual-mask architecture requires simultaneous optimization of depth regression and material domain classification. Conventional light field depth estimation methods optimize a single regression objective, implicitly assuming uniform Lambertian reflectance across the scene. Different from previous works that ignore material-dependent optical variations, our framework explicitly formulates the training objective as a multi-task learning problem. The total loss function $\mathcal{L}_{total}$ integrates a mask-aware depth estimation loss $\mathcal{L}_{depth}$, a domain-aware mask prediction loss $\mathcal{L}_{mask}$, and a physical consistency regularization $\mathcal{L}_{phy}$. This joint optimization forces the network to learn domain-specific feature routing while adhering to the underlying light field physics, ensuring that the depth estimation process is explicitly guided by the material properties of the scene.

\paragraph{Mask-Aware Depth Estimation Loss}
Standard $L_1$ or $L_2$ depth losses treat all pixels equally, making the optimization vulnerable to outliers caused by specular highlights and translucent scattering in non-Lambertian regions. To mitigate this, we design a mask-aware depth loss that applies distinct penalty functions conditioned on the predicted material domain. 

Specifically, for pixels routed to the Lambertian branch, we employ the standard $L_1$ loss, which provides stable gradients for smooth depth surfaces. For pixels routed to the non-Lambertian branch, we utilize the Huber loss to reduce the sensitivity to extreme depth deviations caused by Epipolar Plane Image (EPI) structure corruption. The mask-aware depth loss is defined as:
\begin{equation}
\begin{aligned}
\mathcal{L}_{depth} &= \frac{1}{|\Omega|} \sum_{p \in \Omega} \Big[ \hat{M}_L(p) \ell_1(d(p), \hat{d}(p)) \\
&\quad + \hat{M}_{NL}(p) \ell_{huber}(d(p), \hat{d}(p)) \Big]
\end{aligned}
\end{equation}
where $p$ denotes the spatial pixel index, $\Omega$ is the spatial domain of the central sub-aperture image, $d(p)$ and $\hat{d}(p)$ represent the ground truth and predicted depth values, respectively. $\hat{M}_L(p)$ and $\hat{M}_{NL}(p)$ are the predicted probabilities for the Lambertian and non-Lambertian domains. $\ell_1$ and $\ell_{huber}$ denote the $L_1$ and Huber loss functions. By conditioning the regression penalty on the predicted material mask, the network adaptively adjusts its tolerance to depth outliers, yielding more robust depth maps in mixed-material scenes.

\paragraph{Domain-Aware Mask Prediction Loss}
Accurate pixel-level material classification is a prerequisite for the dual-branch routing mechanism. However, non-Lambertian pixels (e.g., specular reflections, glass) constitute a severe minority class in standard light field datasets, leading to extreme class imbalance. A standard cross-entropy loss would bias the mask predictor toward the majority Lambertian class, degrading the routing accuracy.

To address this, we formulate the mask prediction loss $\mathcal{L}_{mask}$ by combining Focal Loss and Dice Loss. Focal Loss down-weights the contribution of easily classified Lambertian pixels, focusing the optimization on hard non-Lambertian examples. Dice Loss directly optimizes the intersection-over-union of the minority class, preventing the network from collapsing to trivial all-Lambertian predictions.
\begin{equation}
\begin{aligned}
\mathcal{L}_{focal} &= -\frac{1}{|\Omega|} \sum_{p \in \Omega} \Big[ \alpha (1-\hat{M}_{NL}(p))^\gamma \\
&\quad \times M_{NL}(p) \log(\hat{M}_{NL}(p)) + (1-\alpha) \hat{M}_{NL}(p)^\gamma \\
&\quad \times M_L(p) \log(\hat{M}_L(p)) \Big]
\end{aligned}
\end{equation}
\begin{multline}
\mathcal{L}_{dice} = 1 - \Big( 2 \sum_{p \in \Omega} M_{NL}(p)\hat{M}_{NL}(p) \\
+ \epsilon \Big) \Big/ \Big( \sum_{p \in \Omega} M_{NL}(p) \\
+ \sum_{p \in \Omega} \hat{M}_{NL}(p) + \epsilon \Big)
\end{multline}
\begin{equation}
\mathcal{L}_{mask} = \mathcal{L}_{focal} + \lambda_{dice} \mathcal{L}_{dice}
\end{equation}
where $M_L(p)$ and $M_{NL}(p)$ are the ground truth binary masks derived from the three-layer angular signal decomposition. $\alpha$ and $\gamma$ are the balancing factor and focusing parameter for the Focal Loss, respectively. $\epsilon$ is a small constant for numerical stability, and $\lambda_{dice}$ controls the contribution of the Dice term. This composite objective explicitly penalizes the misclassification of non-Lambertian regions, providing reliable prior masks for the subsequent depth estimation branches.

\paragraph{Physical Consistency Regularization}
The core physical limitation of EPI-based methods is the violation of the Angular Cone Constraint on non-Lambertian surfaces, which breaks the linear EPI assumption. While the dual-mask architecture separates the feature processing, the network still requires explicit physical constraints to prevent hallucinating invalid EPI structures in non-Lambertian regions. Different from previous works that rely solely on data-driven regularization, our module explicitly injects the three-layer angular signal physical decomposition ($I(u,v) = DC + a \cdot u + b \cdot v + \varepsilon(u,v)$) into the loss function.

We design two physical consistency terms. First, for the Lambertian domain, the residual $\varepsilon(u,v)$ should approach zero, meaning the angular gradient must strictly align with the predicted depth slope. We define the angular linearity constraint:
\begin{equation}
\begin{aligned}
\mathcal{L}_{linearity} &= \frac{1}{|\Omega|} \sum_{p \in \Omega} \hat{M}_L(p) \\
&\quad \times \left\| \nabla_{\theta} E(p) - \hat{d}(p) \cdot \mathbf{v}_{\theta} \right\|_2^2
\end{aligned}
\end{equation}
where $\nabla_{\theta} E(p)$ is the angular gradient of the EPI patch centered at $p$, and $\mathbf{v}_{\theta}$ is the unit direction vector corresponding to the epipolar line. 

Second, for the non-Lambertian domain, the residual $\varepsilon(u,v)$ captures the high-frequency material reflections. The angular signal must exhibit a high Coefficient of Variation (CV) or angular gradient magnitude, distinguishing it from smooth Lambertian surfaces. We impose a residual regularization to prevent the network from erroneously smoothing non-Lambertian features:
\begin{equation}
\begin{aligned}
\mathcal{L}_{residual} &= \frac{1}{|\Omega|} \sum_{p \in \Omega} \hat{M}_{NL}(p) \\
&\quad \times \max(0, \tau - \text{CV}(p))^2
\end{aligned}
\end{equation}
where $\text{CV}(p)$ is the coefficient of variation of the angular signal at pixel $p$, and $\tau$ is a physical threshold derived from the empirical distribution of non-Lambertian materials. The total physical regularization is $\mathcal{L}_{phy} = \mathcal{L}_{linearity} + \mathcal{L}_{residual}$. By enforcing these physical priors, the network maintains geometric consistency in Lambertian regions while preserving the distinct statistical properties of non-Lambertian residuals.

\paragraph{Adaptive Weighting and Optimization Strategy}
Optimizing a multi-task objective with highly imbalanced domain distributions requires careful weight scheduling and sampling strategies. The total training objective is formulated as:
\begin{equation}
\begin{aligned}
\mathcal{L}_{total} &= \lambda_d \mathcal{L}_{depth} + \lambda_m \mathcal{L}_{mask} \\
&\quad + \lambda_p \mathcal{L}_{phy}
\end{aligned}
\end{equation}
We set the primary weights as $\lambda_d = 1.0$, $\lambda_m = 0.5$, and $\lambda_p = 0.1$. The depth estimation is the primary task, thus receiving the highest weight. The mask prediction acts as an auxiliary routing mechanism, and the physical regularization serves as a soft constraint to guide feature extraction without overpowering the data-driven gradients.

To address the extreme scarcity of non-Lambertian training scenes (only 4 scenes available in the Non-Lambertian Dataset), we implement a Domain-balanced sampling strategy utilizing a WeightedRandomSampler. This sampler assigns higher selection probabilities to batches containing non-Lambertian and Mixed (Urban) domains, ensuring uniform exposure across different material types during each epoch. Combined with the mixed training mechanism that leverages large-scale synthetic data (e.g., 170 scenes from UrbanLF-Syn), this optimization strategy effectively prevents the model from overfitting to the abundant Lambertian data, maintaining robust generalization across diverse light field domains. To empirically validate this cross-domain generalization, we conduct a comprehensive experimental evaluation across multiple representative light field datasets, as detailed in the following section.

\section{Experiments}

\subsection{Datasets}

To comprehensively evaluate the proposed Unified Dual-Mask Physical Model and validate its generalization across varying surface reflectance properties, we construct a multi-domain experimental setup using three representative light field (LF) datasets. These datasets are meticulously selected to cover a wide spectrum of scene types, ranging from ideal Lambertian surfaces to highly complex mixed urban environments and physically challenging non-Lambertian materials. 

\textbf{HCInew Dataset.} The HCInew dataset primarily serves as the benchmark for \textbf{Lambertian} and standard synthetic scenes. It comprises 22 high-quality scenes (14 scenes for training and 8 for testing) (e.g., \textit{boxes}, \textit{cotton}, \textit{dino}) rendered using Blender \cite{dosovitskiy2021}. Each scene consists of $9 \times 9$ (81) angular views with a spatial resolution of $512 \times 512$ pixels. The ground-truth (GT) depth maps are perfectly aligned and generated directly from the 3D rendering engine, providing pixel-accurate depth supervision. The dataset is divided into training and testing sets following the standard stratified split (e.g., \texttt{training/stratified}), ensuring a diverse distribution of geometric structures in both splits. It is chosen to evaluate the model's fundamental capability in extracting Epipolar Plane Image (EPI) structures under ideal physical assumptions.

\textbf{UrbanLF-Syn Dataset.} To assess the model's robustness in \textbf{Mixed/Urban} scenarios, we incorporate the UrbanLF-Syn dataset. This dataset contains 170 synthetic scenes depicting complex urban street views, featuring intricate textures, severe occlusions, and large-scale architectural structures. Similar to HCInew, it provides $9 \times 9$ angular views with high spatial resolution. The GT depth maps are derived from high-fidelity 3D city models. The scenes are partitioned into training and testing subsets. We select this dataset to challenge the model with complex texture variations and real-world structural complexities, where EPI slopes might be locally ambiguous due to fine details and occlusions.

\textbf{Non-Lambertian Dataset.} Evaluating depth estimation on \textbf{Non-Lambertian} surfaces is the core focus of this work. We utilize a specialized Non-Lambertian dataset (derived from the Zhenglong/Wanner HCI collections) that explicitly features specular (highlight), translucent, and scattering surfaces. Each scene contains $9 \times 9$ angular views at a $512 \times 512$ resolution, with GT depth maps generated via synthetic rendering. \textit{Preprocessing and Constraints:} Originally, the HCI-Old dataset was considered; however, it was excluded from our pipeline due to inconsistent data formatting and lack of standardized depth calibration parameters. Consequently, the available Non-Lambertian dataset is severely constrained, containing only 4 training scenes. Despite this extreme data scarcity, it is indispensable for our study as it introduces the ultimate physical challenge: the violation of the standard EPI linearity assumption caused by view-dependent reflectance (e.g., specular highlights shifting across views).

\textbf{Data Preprocessing and Domain-Balanced Sampling.} 
A critical challenge in our multi-domain setup is the severe data imbalance across different scene types (e.g., 170 scenes in UrbanLF-Syn versus only 4 scenes in the Non-Lambertian dataset). To prevent the network from overfitting to the data-rich domains and neglecting the physically challenging non-Lambertian cases, we implement a rigorous \textbf{domain-balanced sampling strategy}. Specifically, we employ a \texttt{WeightedRandomSampler} during the training phase. Each sample is assigned a weight inversely proportional to the total number of samples in its respective domain (Lambertian, Mixed/Urban, or Non-Lambertian). This ensures that each mini-batch maintains an equitable exposure to all three domains, facilitating stable multi-task learning and joint optimization of the depth and domain classification losses.

The key characteristics and configurations of the datasets utilized in our experiments are summarized in Table~\ref{tab:datasets}.

\begin{table*}[!t]
\caption{Summary of Datasets Used for Training and Evaluation}
\label{tab:datasets}
\centering\small
\resizebox{\textwidth}{!}{%
\begin{tabular}{lllllll}
\toprule 
\textbf{Dataset} & \textbf{Domain / Scene Type} & \textbf{No. of Scenes (Train / Test)} & \textbf{Angular Views} & \textbf{Spatial Resolution} & \textbf{GT Depth Source} & \textbf{Key Characteristics \& Challenges} \\
\midrule 
\textbf{HCInew} & Lambertian & 22 (14 / 8) & $9 \times 9$ (81) & $512 \times 512$ & Synthetic (Blender) & Ideal EPI structures; standard geometric shapes and textures. \\
\textbf{UrbanLF-Syn} & Mixed / Urban & 170 (Partitioned) & $9 \times 9$ (81) & $512 \times 512$ & Synthetic (3D City Models) & Complex textures, severe occlusions, large-scale urban structures. \\
\textbf{Non-Lambertian} & Non-Lambertian & 4 / (Cross-validated) & $9 \times 9$ (81) & $512 \times 512$ & Synthetic (Blender) & Specular, translucent, scattering surfaces; EPI physical assumption failure; extreme data scarcity. \\
\bottomrule
\end{tabular}}
\vspace{4pt}
\parbox{\textwidth}{\footnotesize \textit{Note: The HCI-Old dataset was excluded during preprocessing due to inconsistent data formatting and lack of standardized depth calibration parameters. Consequently, the Non-Lambertian training set is limited to 4 scenes.}}
\end{table*}

\subsection{Implementation Details}

\textbf{Hardware and Software Environment.} 
The proposed Unified Dual-Mask Physical Model is implemented using the PyTorch framework. All experiments are conducted on a high-performance workstation equipped with four NVIDIA RTX 3090 GPUs (24GB VRAM each) and an AMD EPYC 7742 CPU. The Distributed Data Parallel (DDP) strategy is utilized to synchronize gradients and accelerate the training process across multiple GPUs.

\textbf{Training Hyperparameters.} 
We employ the AdamW optimizer to train the network, initialized with a learning rate of $1 \times 10^{-4}$ and a weight decay of $1 \times 10^{-4}$. The learning rate is dynamically adjusted using a Cosine Annealing scheduler, decaying to a minimum of $1 \times 10^{-6}$ over the course of training. The batch size is set to 32 per GPU, and gradient accumulation is applied with 2 steps, yielding an effective global batch size of 256. The model is trained for a total of 150 epochs. To mitigate the severe data imbalance among the Lambertian, Mixed/Urban, and Non-Lambertian datasets (notably, the Non-Lambertian dataset contains only 4 training scenes), we adopt a domain-balanced sampling strategy utilizing a \texttt{WeightedRandomSampler}. This ensures equitable exposure to each physical domain during every epoch. The key hyperparameters are summarized in Table~\ref{tab:hyperparameters}.

\begin{table}[!t]
\caption{Summary of Key Training Hyperparameters}
\label{tab:hyperparameters}
\centering
\resizebox{\columnwidth}{!}{%
\begin{tabular}{ll}
\toprule 
\textbf{Hyperparameter} & \textbf{Value} \\
\midrule 
Optimizer & AdamW \\
Initial Learning Rate & $1 \times 10^{-4}$ \\
LR Scheduler & Cosine Annealing \\
Minimum Learning Rate & $1 \times 10^{-6}$ \\
Batch Size (per GPU) & 32 \\
Gradient Accumulation Steps & 2 \\
Effective Batch Size & 256 \\
Total Epochs & 150 \\
Weight Decay & $1 \times 10^{-4}$ \\
Sampler & WeightedRandomSampler (Domain-balanced) \\
\bottomrule
\end{tabular}}
\end{table}

\textbf{Data Augmentation.} 
To improve the generalization capability and prevent overfitting---particularly given the extreme scarcity of Non-Lambertian training samples---we apply a comprehensive set of data augmentation techniques tailored for 4D light field data. These include random spatial cropping (resizing to $512 \times 512$ resolution), random horizontal and vertical flips (applied consistently across all angular views to preserve epipolar geometry), and random angular shifts. Additionally, we apply color jittering (random adjustments to brightness, contrast, and saturation) and Gaussian noise injection. These photometric augmentations simulate sensor noise and varying illumination conditions, which is highly beneficial for robust material recognition and domain classification.

\textbf{Loss Function Weights.} 
The network is optimized in a multi-task learning paradigm, jointly minimizing the depth estimation loss and the auxiliary domain classification loss. The overall objective function is defined as $\mathcal{L}_{total} = \lambda_{depth} \mathcal{L}_{depth} + \lambda_{domain} \mathcal{L}_{domain}$. We empirically set the loss weights to $\lambda_{depth} = 1.0$ and $\lambda_{domain} = 0.1$ to prioritize accurate depth regression while providing sufficient gradient signals for domain-aware feature disentanglement.

\subsection{Evaluation Metrics}

To quantitatively assess the performance of the proposed Unified Dual-Mask Physical Model and ensure fair comparisons with existing state-of-the-art methods, we employ a comprehensive set of evaluation metrics. Following standard protocols, we use Mean Squared Error (MSE) and Bad Pixel ratio as our primary evaluation metrics for depth estimation accuracy.

Specifically, the \textbf{Mean Squared Error (MSE)} measures the average squared difference between the estimated depth map and the ground-truth depth map, providing a global assessment of regression accuracy. Lower MSE values indicate better overall depth estimation performance.

Furthermore, to evaluate local precision and the proportion of highly inaccurate predictions, we compute the \textbf{Bad Pixel ratio}. A pixel is classified as a ``bad pixel'' if the absolute error between the predicted depth and the ground truth exceeds a predefined threshold. Following standard protocols in light field depth estimation, we report the Bad Pixel ratios at two strict thresholds: $<0.01$ and $<0.03$ (normalized by the maximum depth range). A lower Bad Pixel ratio signifies that the model produces fewer severe outliers, which is particularly crucial for challenging non-Lambertian regions where EPI structures are heavily distorted.

Additionally, we report the \textbf{Mean Absolute Error (MAE)} to provide a robust measure of the average magnitude of errors, complementing the MSE metric. All metrics are computed exclusively on the valid regions of the depth maps, ignoring masked or invalid background pixels.

\section{Conclusion}

In this paper, we present a unified dual-mask physical model for light field depth estimation that explicitly addresses the breakdown of the Epipolar Plane Image (EPI) assumption in non-Lambertian environments. By integrating physical signal decomposition with a dual-mask routing mechanism, the proposed framework effectively decouples depth and material information under low angular resolution conditions. Extensive evaluations demonstrate that our method achieves an overall Mean Absolute Error (MAE) of 0.133 and an MAE of 0.081 in mixed urban scenes, validating the efficacy of the domain-balanced sampling strategy in handling long-tailed data distributions and improving generalization capabilities.

The core contributions of this work are summarized as follows:
\begin{itemize}
\item \textbf{Three-layer angular signal physical decomposition and geometric material feature extraction:} This mechanism decouples depth and material information from physical and geometric perspectives, overcoming the bottleneck of traditional frequency-domain methods in reliably classifying materials under discrete sampling. It provides a robust feature foundation for the precise identification of non-Lambertian regions, preventing the network from blindly learning noise that violates physical assumptions.
\item \textbf{Unified depth estimation architecture with dual-mask routing and domain-balanced sampling:} The dual-mask mechanism mitigates the error amplification caused by the global enforcement of a single EPI assumption, enhancing robustness on complex reflective surfaces. Concurrently, the domain-balanced sampling strategy exploits the training potential of large-scale data under long-tailed distributions, significantly improving the overall depth estimation accuracy in mixed and urban scenes.
\item \textbf{Quantitative analysis and theoretical verification system for the physical root causes of EPI assumption failure:} This system redefines the failure of light field depth estimation in non-Lambertian domains as the destruction of underlying physical assumptions rather than improper network parameter or loss function design. It provides a solid theoretical basis for transitioning away from pure EPI-based methods, effectively avoiding invalid computational consumption in incorrect directions and enhancing the trial-and-error efficiency of automated research systems.
\end{itemize}

The theoretical verification system established in this work redirects the focus of light field depth estimation from purely geometric EPI structures to physically aware modeling. These findings facilitate the development of next-generation automated visual systems operating in complex, real-world reflective environments.

While the proposed three-layer angular signal decomposition, formulated as $I(u,v) = DC + a\cdot u + b\cdot v + \epsilon(u,v)$, effectively isolates material residuals from ambient illumination and linear disparity, it inherently assumes locally linear epipolar geometry. In scenarios characterized by severe anisotropic reflections or complex interreflections---such as translucent subsurface scattering or concave metallic surfaces---the linear approximation $a\cdot u + b\cdot v$ may underfit the true disparity field. Consequently, residual depth cues can leak into the material term $\epsilon(u,v)$, potentially degrading the reliability of the geometric classifier. This highlights a fundamental physical trade-off: maintaining model parsimony for robust estimation under low angular resolution (e.g., $9 \times 9$) versus capturing higher-order bidirectional reflectance distribution function (BRDF) dynamics.

From a systems and computational perspective, replacing traditional frequency-domain analyses with geometric angular features introduces critical advantages for real-time processing. Frequency-domain transformations on sparse angular grids suffer from severe spectral leakage and necessitate extensive zero-padding, which drastically inflates memory bandwidth and computational latency. Conversely, our geometric feature extraction---relying on angular gradients and correlation coherence---operates directly in the spatial-angular domain with strictly $\mathcal{O}(N)$ complexity and highly localized memory access patterns. This makes the material prior generation exceptionally well-suited for edge deployment and light field video pipelines. However, this discrete geometric approximation may discard high-frequency angular modulations that dense sampling regimes (e.g., $>17 \times 17$) could otherwise resolve, suggesting a potential accuracy ceiling for extremely fine specularities.

Furthermore, the dual-mask routing mechanism provides benefits that extend beyond explicit feature separation; it fundamentally resolves the optimization conflicts inherent in mixed-material scenes. Lambertian and non-Lambertian regions exhibit drastically different Epipolar Plane Image (EPI) gradient distributions. Forcing a unified backbone to regress both simultaneously induces severe negative transfer and gradient interference during backpropagation. By decoupling the gradient flows via mask-guided routing, the network prevents the averaging of conflicting loss signals. Future iterations could replace hard binary masks with soft, confidence-weighted routing derived from the classifier's probabilistic outputs, thereby enabling smoother gradient transitions at material boundaries and reducing edge artifacts in the final depth maps.

Looking forward, extending this physical decomposition to temporal light field video presents a compelling direction, as motion parallax introduces complex coupling between the $DC$ and linear terms. Additionally, integrating our geometric material priors with modern neural rendering frameworks, such as 3D Gaussian Splatting, could facilitate physically plausible novel view synthesis in highly specular and non-Lambertian environments.

Despite the demonstrated efficacy of the proposed unified dual-mask physical model, several limitations remain that outline promising directions for future research. 

First, the current three-layer angular signal decomposition relies on a first-order linear approximation for the disparity term ($a\cdot u + b\cdot v$) and a constant ambient illumination assumption (the DC term). While this formulation effectively decouples depth and material under standard conditions, it struggles with complex global illumination effects, such as inter-reflections and ambient occlusion, as well as highly curved surfaces where the epipolar geometry exhibits higher-order non-linearities. In such scenarios, the residual term $\epsilon(u,v)$ inadvertently absorbs geometric distortions, potentially degrading the material classifier's accuracy. Furthermore, the geometric angular features are primarily tailored for opaque specular and diffuse surfaces. For translucent materials or subsurface scattering, the angular radiance distribution follows complex volumetric transport equations, making the current handcrafted geometric priors less discriminative. Moreover, although the feature extraction is robust at $9 \times 9$ angular resolution, its stability fundamentally degrades in extremely sparse setups (e.g., $2 \times 2$), where the statistical variance of the correlation coherence becomes prohibitively high.

To address these limitations, future work will explore integrating higher-order physical representations. Specifically, replacing the explicit polynomial decomposition with implicit neural rendering frameworks, such as Neural Radiance Fields (NeRF) or 3D Gaussian Splatting (3DGS), could natively model complex light transport and view-dependent material properties without relying on strict linear Epipolar Plane Image assumptions. Additionally, we plan to extend the dual-mask routing architecture to dynamic light field videos (4D light fields). This extension requires incorporating temporal consistency constraints to handle non-rigid deformations and time-varying non-Lambertian phenomena, such as fluid dynamics or moving specular highlights. Finally, to alleviate the reliance on synthetic datasets with perfect material annotations, we will investigate self-supervised material-depth disentanglement strategies. By leveraging physical consistency losses across multi-view photometric stereo setups, the network could learn to intrinsically separate reflectance and geometry directly from raw, unannotated real-world light fields, thereby bridging the persistent domain gap between synthetic training environments and in-the-wild captures.

\begin{figure*}[!t]
\centering
\includegraphics[width=\textwidth]{./output/figures/fig1.pdf}
\caption{The overall framework of the proposed method. It explicitly decouples depth and material properties via a three-layer angular signal decomposition and routes features into Lambertian and non-Lambertian branches guided by a per-pixel geometric mask.}
\label{fig:fig1}
\end{figure*}

\begin{figure}[!t]
\centering
\includegraphics[width=\columnwidth]{./output/figures/fig3.pdf}
\caption{The per-pixel geometric mask generator utilizes angular gradients and coherence features to classify Lambertian and non-Lambertian regions via a Random Forest, providing reliable prior masks for the dual-branch routing.}
\label{fig:fig3}
\end{figure}

\begin{figure*}[!t]
\centering
\includegraphics[width=\textwidth]{./output/figures/fig4.pdf}
\caption{Quantitative and qualitative analysis of EPI assumption failures on non-Lambertian surfaces. Specular and scattering surfaces violate the angular cone constraint, resulting in broken EPI lines and X-shaped bimodal distributions.}
\label{fig:fig4}
\end{figure*}

\begin{figure}[!t]
\centering

\includegraphics[width=\columnwidth]{./output/figures/fig5.pdf}
\caption{The progressive training strategy with domain-balanced sampling mitigates data scarcity in non-Lambertian domains by balancing synthetic and real-world scene distributions to prevent severe overfitting.}
\label{fig:fig5}
\end{figure}

\begin{figure*}[!t]
\centering
\includegraphics[width=\textwidth]{./output/figures/fig6.pdf}
\caption{Visual comparison of depth estimation results on complex mixed/urban scenes. Our method accurately recovers depth on highly reflective and scattering non-Lambertian surfaces where state-of-the-art baselines severely degrade.}
\label{fig:fig6}
\end{figure*}



\bibliographystyle{IEEEtran}
\bibliography{references}

\end{document}
